\section{Rezultate}
\label{sec:results}

Proiectul a fost testat în Unity 6000.4.2f1, pe un circuit construit prin sistemul de waypoint-uri descris în secțiunea de metodologie.

\subsection{Evaluarea Fizicii Vehiculului}

Parametrii vehiculului jucătorului și efectele lor sunt sumarizate în Tabelul~\ref{tab:vehicle_params}.

\begin{table}[h]
\centering
\small
\begin{tabular}{lrl}
\toprule
\textbf{Parametru} & \textbf{Valoare} & \textbf{Efect} \\
\midrule
Masă vehicul & 1500 kg & Stabilitate \\
Cuplu maxim motor & 8000 Nm & Accelerare \\
Viteză maximă & 320 km/h & Limita superioară \\
Forță de frânare & 15000 Nm & Frânare eficientă \\
Coeficient downforce & 50 & Aderență la viteză \\
Asistență laterală & 5 & Anti-derapaj \\
Fricțiune normală & 3.5 & Grip anvelope \\
\bottomrule
\end{tabular}
\caption{Parametrii fizici ai vehiculului jucătorului.}
\label{tab:vehicle_params}
\end{table}

Testele au demonstrat:
\begin{itemize}
    \item \textbf{Accelerare progresivă:} Factorul de boost ($2.5\times$) la viteze mici asigură o pornire rapidă.
    \item \textbf{Stabilitate la viteză mare:} Downforce-ul și asistența laterală elimină oscilațiile.
    \item \textbf{Controlul tracțiunii:} Limita de 15 km/h previne patinarea roților.
    \item \textbf{Coasting natural:} Drag-ul aerodinamic oferă decelerare graduală.
\end{itemize}

\subsection{Evaluarea Inteligenței Artificiale}

\begin{table}[h]
\centering
\small
\begin{tabular}{ll}
\toprule
\textbf{Criteriu} & \textbf{Rezultat} \\
\midrule
Navigare completă traseu & Da \\
Evitarea blocajelor (respawn) & Activ (10s) \\
Variabilitate traiectorie & Aleatorizare pe lățime \\
Adaptare zone boost/breakers & Funcțională \\
\bottomrule
\end{tabular}
\caption{Evaluarea calitativă a sistemului AI.}
\label{tab:ai_eval}
\end{table}

\subsection{Definirea Traseului prin Waypoint-uri}

Sistemul de waypoint-uri permite construirea unor trasee de forme arbitrare direct din editorul Unity, cu vizualizare în Scene View a conexiunilor și a lățimii fiecărui punct. Aceeași secvență de waypoint-uri este folosită atât pentru poziționarea segmentelor de drum și a pereților, cât și pentru navigarea AI, eliminând duplicarea datelor între reprezentarea vizuală și logica de navigare.

\subsection{Sistemele de Joc}

Numărătoarea inversă, sistemul de tururi, meniul de pauză și ecranul de final funcționează conform așteptărilor, cu gestionare corectă a \texttt{Time.timeScale} și a cursorului mouse.
